\documentclass[12pt]{ximera}
\title{Homework 5: Banzhaf Power in the Boone City Council}
\begin{document}
\begin{abstract}
A real weighted-voting system analyzed with coalition ``types'': the Boone city
council, where the mayor holds zero votes but a veto that a two-thirds majority can
override. The Banzhaf power distribution turns out to be a surprise.
\end{abstract}
\maketitle

\section*{Banzhaf Power with Coalition Types}

The Boone city council consists of the mayor ($M$) and seven council members
($C_1,\ldots,C_7$). Since the council members have indistinguishable roles, we adopt
the shorthand ``$C$'' to stand for any one council member. Each council member has one
vote.

Per the state of Iowa's city council governing documents, the mayor has \emph{zero}
votes. Instead, the mayor has a veto, which can be overruled by a two-thirds majority
among the seven council members. For simplicity, we include the mayor in a coalition
whenever he does \emph{not} intend to veto, and omit him whenever he does.

A simple majority among the seven council members is needed to pass a motion, and a
two-thirds majority of council members can overrule a veto.

\begin{problem}
\begin{prompt}
\textbf{(a)} How many council members does a coalition need to pass a motion
\emph{with} mayoral support?

$\answer{4}$

\begin{solution}
A simple majority of $7$ council members means more than $3.5$, so at least $4$.
\end{solution}
\end{prompt}

\begin{prompt}
\textbf{(b)} How many council members does a coalition need to pass a motion
\emph{without} mayoral support?

$\answer{5}$

\begin{hint}
Without mayoral support the motion is vetoed, so the council must override. Two-thirds
of $7$ is not a whole number --- round up.
\end{hint}

\begin{solution}
Overriding the veto takes a two-thirds majority of the $7$ council members, and
\[ \tfrac23 \cdot 7 = 4.67, \]
so at least $5$ council members are required.
\end{solution}
\end{prompt}

\begin{prompt}
\textbf{(c)} Below is every coalition type. Identify which are winning, and how many
players are critical in each.

For coalitions \emph{with} mayoral support:

$\set{M, C, C, C, C}$ (4 council members, $35$ distinct coalitions) --- is it winning?
\begin{multipleChoice}
\choice[correct]{Yes}
\choice{No}
\end{multipleChoice}
Number of critical players in each such coalition: $\answer{5}$

\smallskip
$\set{M, C, C, C, C, C}$ (5 council members, $21$ coalitions) --- number of critical
players: $\answer{0}$

\smallskip
For coalitions \emph{without} mayoral support:

$\set{C, C, C, C, C}$ (5 council members, $21$ coalitions) --- number of critical
players: $\answer{5}$

\smallskip
$\set{C,C,C,C,C,C}$ (6 council members, $7$ coalitions) --- number of critical
players: $\answer{0}$

\begin{hint}
A player is critical when removing them turns a winning coalition into a losing one.
Careful with $\set{M,C,C,C,C,C}$: if the mayor walks away, the five remaining council
members can still override his veto.
\end{hint}

\begin{solution}
The winning types are: with the mayor, $4$ or more council members; without the mayor,
$5$ or more.

\textbf{$\set{M,C,C,C,C}$ --- winning, $5$ critical players.} Remove $M$ and the four
remaining council members fall short of the $5$ needed to override, so $M$ is
critical. Remove any one $C$ and only $3$ council members remain with mayoral support,
short of $4$ --- so all four council members are critical too. That is $1+4 = 5$
critical players.

\textbf{$\set{M,C,C,C,C,C}$ --- winning, but nobody is critical.} Remove $M$ and the
$5$ council members override the veto and still win. Remove one $C$ and you are left
with $\set{M,C,C,C,C}$, which also still wins. The same holds for every larger
coalition with the mayor.

\textbf{$\set{C,C,C,C,C}$ --- winning, all $5$ critical.} Removing any member leaves
$4$, which cannot override the veto.

\textbf{$\set{C,C,C,C,C,C}$ --- winning, nobody critical.} Removing one still leaves
$5$, enough to override.

Every other type is either losing (too few council members) or winning with no
critical players.
\end{solution}
\end{prompt}

\begin{prompt}
\textbf{(d)} What is the total critical count in this system?

$\answer{280}$

\begin{hint}
Only two coalition types contribute. Multiply the number of critical players per
coalition by how many distinct coalitions there are of that type.
\end{hint}

\begin{solution}
Only two types contribute anything:
\[
\underbrace{35 \times 5}_{\set{M,C,C,C,C}} + \underbrace{21 \times 5}_{\set{C,C,C,C,C}}
= 175 + 105 = 280.
\]
\end{solution}
\end{prompt}

\begin{prompt}
\textbf{(e)} Calculate the mayor's critical count.

$\answer{35}$

\begin{solution}
The mayor is critical only in the type $\set{M,C,C,C,C}$, and there are
$C(7,4)=35$ such coalitions. So $M$'s critical count is $35$.
\end{solution}
\end{prompt}

\begin{prompt}
\textbf{(f)} Calculate the critical count of the council members as a group, and then
for each individual council member.

As a group: $\answer{245}$ \qquad Each individual council member: $\answer{35}$

\begin{solution}
The group total is everything left over: $280 - 35 = 245$.

Directly: council members are critical $35\times 4 = 140$ times in the
$\set{M,C,C,C,C}$ type and $21\times 5 = 105$ times in the $\set{C,C,C,C,C}$ type, and
$140+105 = 245$. \checkmark

By symmetry the seven members share this equally: $245/7 = 35$ each.
\end{solution}
\end{prompt}

\begin{prompt}
\textbf{(g)} Give the Banzhaf Power Distribution of this system across all eight
players, as percentages.

Mayor: $\answer[tolerance=0.05]{12.5}\%$ \qquad
Each council member: $\answer[tolerance=0.05]{12.5}\%$

\begin{solution}
\[
\text{Mayor}: \frac{35}{280}=\frac18=12.5\%, \qquad
\text{Each council member}: \frac{35}{280}=12.5\%.
\]
All eight players have \emph{exactly the same} Banzhaf power, and the eight shares
total $100\%$. \checkmark

This is the surprise. The mayor has zero votes and cannot vote on anything --- yet his
veto gives him precisely as much power as a full voting member of the council. Raw
vote counts can be a very poor guide to actual influence.
\end{solution}
\end{prompt}

\begin{prompt}
\textbf{(h)} What is one way to write this eight-player system using bracket notation?

$[\answer{5} : 1,1,1,1,1,1,1,1]$

\begin{hint}
Part (g) says all eight players have equal power, which suggests giving them all equal
weight. How many of the eight must agree for a motion to pass?
\end{hint}

\begin{solution}
Give every player, including the mayor, one vote, and set the quota at $5$:
\[ [5:1,1,1,1,1,1,1,1]. \]

Check that this reproduces the original rules exactly:
\begin{itemize}
\item Mayor plus $4$ council members $= 5$ players, which wins --- matching part (a).
\item Mayor plus $3$ council members $= 4$, which loses. \checkmark
\item $5$ council members without the mayor $= 5$, which wins --- matching part (b).
\item $4$ council members without the mayor $= 4$, which loses. \checkmark
\end{itemize}
The winning coalitions are identical, so the systems are equivalent --- and in this
form the equal power distribution from part (g) is obvious by symmetry.
\end{solution}
\end{prompt}
\end{problem}

\end{document}
